\documentclass[11pt]{article}
\usepackage[T1]{fontenc}
\usepackage{textcomp}
\usepackage[margin=0.75in]{geometry}
\usepackage{amsmath,amssymb,amsthm}
\usepackage{array,booktabs}
\usepackage{hyperref}
\setlength{\parindent}{0pt}
\newtheorem{theorem}{Theorem}
\theoremstyle{definition}
\newtheorem{definition}{Definition}
\title{Flow Proof Artifact: Nat-mul-fifteen-one}
\author{Generated by \texttt{flow doc proof}}
\date{Flow Proof Book}
\begin{document}
\maketitle
Fifteen times one is fifteen.\\[0.5em]
\textbf{Source.} peano — https://en.wikipedia.org/wiki/Peano\_axioms\\[0.5em]
\bigskip
\noindent{\large \textbf{Derived fact 1.} \textit{15 * 1 = 15} \hfill the natural numbers \cdot multiplication \cdot fifteen times one is fifteen}\par
\smallskip\noindent\textit{Coordinate: the natural numbers \cdot multiplication \cdot fifteen times one is fifteen}\par
\smallskip\noindent\textit{Source: peano}\par
\smallskip\noindent\textit{Built on: fourteen times one is fourteen, for multiplication on the natural numbers}\par
\medskip
\noindent\textbf{Goal.} 15 * 1 = 15\par
\noindent$\displaystyle 15 \cdot 1 = 15$\par
\medskip
\noindent
\begin{tabular}{@{} >{\raggedright\arraybackslash}p{0.06\textwidth} >{\raggedright\arraybackslash}p{0.47\textwidth} >{\raggedleft\arraybackslash}p{0.41\textwidth} @{}}
\textbf{\#} & \textbf{Proof} & \textbf{Mathematics} \\
\midrule
\textbf{1.} & We prove that fifteen times one is fifteen for multiplication on the natural numbers. & \\[0.45em]
\textbf{2.} & We invoke the derived fact governing multiplication on the natural numbers: fourteen times one is fourteen, for multiplication on the natural numbers. & \\[0.45em]
\textbf{3.} & From step 2, this implies 15 times 1 equals 15. Hence proven. & $\displaystyle 15 \cdot 1 = 15$ \\[0.45em]
\bottomrule
\end{tabular}
\label{thm:Nat-multiplication-fifteen_times_one_is_fifteen}
\par\medskip\hrule
\end{document}
